% =====================================================================
%  8 Discussion
%  Sources: reports/v3_runC_results.md 5-6, reports/v3_identity_vs_baseline.md,
%  CLAUDE.md section 2 (story), the "honest limits" in
%  reports/proposal_sketch.md section 7.
% =====================================================================
\section{Discussion}
\label{sec:discussion}

\subsection{A model-made group effect}
% Answer RQ1 and RQ2 together: the identity changes the answer (region),
% partly via the place name; the change is not what real groups do.
% Framing moves models away from every group; Llama reproduces the
% direction of regional offsets but exaggerates them (stereotype-like?).
\fillin{Interpretation of the behavioural results.}

\subsection{Calibration is not alignment}
% Why the supervisor's criterion can be met by a constant shift; what a
% better criterion would need (ranking, relative alignment, flag rate,
% soft-label metrics). The Llama dampener heads as the constructive
% result: "removes an identity bias", not "thinks like the group".
\fillin{Interpretation of the steering results.}

\subsection{What the data can and cannot support}
% No reliable item-level group signal in D3CODE (demographic or moral);
% disagreement is mostly within groups. So "make the model act like
% group X" has no item-level target here. Links: Hu & Collier (variance
% explained < 10%), perspectivism literature.
\fillin{Implications of the group-signal analysis.}

\subsection{Implications for LLM-based moderation}
% Over-flagging (recall 0.99, precision 0.26) as the practical issue;
% persona prompts make it worse; threshold / calibration matters more
% than personas; borderline texts.
\fillin{Practical implications.}

\subsection{Lessons for interpretability}
% Attention patterns are a poor guide to causal heads (overlap barely
% above chance); item-level significance under strong single-head
% steering is meaningless without matched nulls; information reaches
% the answer via intermediate tokens.
\fillin{Methodological lessons.}

\subsection{Limitations}
\label{sec:limitations}
% - One dataset, English only, two 7-8B open models, one prompt family,
%   single 0/1 answer token.
% - 500 borderline texts (band filter): results describe borderline
%   texts, not easy ones.
% - Small rater groups per text (>= 4): noisy group shares.
% - Only 2 phrases for gender, religion and identity: low power.
% - Steering strength 100 is extreme; batched bf16 noise up to ~0.5
%   log-odds (handled by within-batch comparisons).
% - Constant shift needs labelled data too.
% - Place controls and steering only for 4 of 8 regions.
\fillin{Limitations.}

\subsection{Ethical considerations}
% Persona prompting can reproduce stereotypes (Llama's exaggerated
% regional offsets); risks of "simulating" groups; D3CODE consent and
% licence; offensive example texts in the thesis (content warning?).
\fillin{Ethics statement.}
